\section*{ID7400: Additional Effects That Reduce Light Intensity}
\paragraph{Metadata.}
PiID: 2276930624743 | RTEID: RTE:P28735.8434:T3.3655 | UUID: 61d055fc-7841-5bfe-ab2f-1b31361a1b48

\paragraph{Canonical Statement.}
Intensity received by a telescope: \textbackslash{}[ I = \textbackslash{}frac\{I\_0 \textbackslash{}cdot A \textbackslash{}cdot \textbackslash{}text\{albedo\}\}\{4 \textbackslash{}pi d\textasciicircum{}2\} \textbackslash{}] - \textbackslash{}( I\_0 \textbackslash{}) = incident solar intensity - \textbackslash{}( A \textbackslash{}) = cross-sectional area of the planet - \textbackslash{}( d \textbackslash{}) = distance from planet to observer

\paragraph{Canonical Equation.}
\[
I = \frac{I_0 \cdot A \cdot \text{albedo}}{4 \pi d^2}
\]

\paragraph{Dictionary.}
Additional Effects That Reduce Light Intensity — I = (I\_0 · A · albedo)/(4 π d\textasciicircum{}2)

\paragraph{Encyclopedia.}
Additional Effects That Reduce Light Intensity is an algebraic definition, written I = (I\_0 · A · albedo)/(4 π d\textasciicircum{}2). It is an algebraic definition expressing the quantity directly in terms of the variables on its right-hand side. It is built from π, I\_0, A, defining I. Within the Trawin Zero Point registry it serves as one element of the framework's mathematical narrative.
